% =====================================================================
%  Section 5, Discussion and Limitations (English)
%
%  "What this study does not claim" belongs here, not in Section 4 (Results).
%  Length target: ~0.6 pp.
% =====================================================================

\section{Discussion and Limitations}
\label{sec:discussion}

What SPR-IRL recovers is a picture that places the preferences of three investor
types side by side on one common scale. Foreign and retail investors diverge in
opposite directions on momentum, the same contrast reappears in the context
response, and institutional investors show no identified preference once herding
is set aside. All three findings come from giving every type the same features
and the same model form and estimating only the parameters separately, so the
contrast was not imposed by design.

\subsection{What This Study Does Not Claim}
\label{sec:notclaim}

As the final stage of the protocol we state what the results do not support.

\textbf{(1) Return predictability.} Our metrics are reproduction errors for the
direction and intensity of order flow and include no returns, Sharpe ratios or
transaction costs. No trading rule follows from Table~\ref{tab:performance}.

\textbf{(2) Superior flow prediction.} The persistence baseline beats SPR-IRL on
directional accuracy and correlation (\S\ref{sec:oos}). Preference recovery and
prediction call for different feature sets. In a preliminary check, adding the
lagged own action as a fourth feature lifts the model above the baseline
(directional accuracy $0.646/0.569/0.629$) but reverses the sign of the retail
$\alpha^{KOSPI}$ and eliminates the underwater coefficient. \textbf{Some of the
mirror images reported here are therefore not separable from the autocorrelation
of order flow.}

\textbf{(3) Causal effects.} The recovered weights are conditional associations
among standardized observed features. Because feature--context correlations range
from $-0.29$ to $0.22$, the magnitude of any coefficient depends on the other
terms being controlled for.

\textbf{(4) A behavioural reading of the negative cross-type association.} The
herding coefficient and the mirrored context responses are not separable from the
market-clearing constraint (\S\ref{sec:beta}).

\textbf{(5) A test of the disposition effect.} The retail underwater sign agrees
with the literature but does not pass the month-block bootstrap, and the
average-cost proxy is built from aggregate trades rather than account-level
realized gains.

\textbf{(6) Generalization across stocks.} The single-stock design is a control
for isolating heterogeneity across types; extension to other stocks and markets
is untested.

\subsection{Future Work}

The most pressing task is to add the lagged own action as a feature and check
whether the mirror images survive once autocorrelation is controlled. Next,
decomposing the ``institutional'' category into pension funds, investment trusts,
insurers and others may reveal preferences currently hidden by offsetting flows.
Finally, the myopia assumption remains to be relaxed. As
Proposition~\ref{prop:lasso} shows, a one-step-ahead setting makes the reward
coincide with the policy; introducing state dynamics and multi-period value
separates the two and turns the exercise into a structural recovery that
regression cannot reach. The symmetric estimation and validation protocol of
SPR-IRL carry over unchanged to that extension.
